% Chapter 6 -- Taking Care of Your Work.  Sources (Law 1): os/milton/kmain.c
% (Shut Down, Restart and Exit to InitechDOS commit DESKTOP.DB and ask the
% open application to quit; nothing else writes the desktop arrangement on
% a power loss); os/flair/finder_ops.c (the Trash moves, Empty Trash
% deletes); os/apps/itext/itext.c (an unsaved document lives in memory
% only); the InitechDOS books for COPY and XCOPY.
% GATES: test-flair-shutdown test-flair-empty-trash test-flair-itext test-desktop-db
\Gates{test-flair-shutdown, test-flair-empty-trash, test-flair-itext,
  test-desktop-db}
\chapter{Taking Care of Your Work}

A few habits protect the documents you make with FLAIR.

\section{Save often}
\index{saving}
A document you are working on lives in the computer's memory until you save
it.  If the computer is turned off, or an application stops working, what
you have not saved is lost.  Save with \key{Ctrl-S} every few minutes, and
always before you leave your desk.

\section{Shut down before you turn off}
\index{shutting down}
Choose \menuitem{Shut Down} from the Finder's \menuname{Special} menu and
wait for the message that it is safe to switch off.  Shut Down asks the
open application to save your work and records the arrangement of your
desktop on the data disk.  Turning the computer off without it loses both.

\section{Look in the Trash before you empty it}
\index{Trash}
What you drag to the Trash is only put aside.  \menuitem{Empty Trash}
deletes it for good.  Double-click the Trash and check its window before
you empty it.

\section{Keep copies of your disks}
\index{backup copies}
Keep a copy of your data disk, and bring it up to date at the end of each
day.  Leave FLAIR with \menuitem{Exit to InitechDOS} and copy your files
with the \file{COPY} or \file{XCOPY} command; the \doctitle{InitechDOS
User's Guide} shows you how.  Store the copies away from the computer.

\section{Handling disks}
\index{disks, handling}
\begin{itemize}
\item Keep disks away from magnets, telephones, displays and heat.
\item Do not touch the surface of a disk, and do not bend it.
\item Remove a disk from its drive only when the computer is turned off.
\item Label each disk, and write on the label before you put it on the
  disk.
\end{itemize}
